\documentclass[a4paper,10pt]{article}
\usepackage[margin=14mm]{geometry}
\usepackage{fontspec}
\usepackage[french]{babel}
\usepackage{xcolor}
\usepackage{enumitem}
\usepackage{titlesec}
\usepackage{hyperref}
\definecolor{accent}{HTML}{176B70}
\definecolor{muted}{HTML}{52616B}
\hypersetup{colorlinks=true,urlcolor=accent,pdftitle={Léo Martin - Product Builder Engineer},pdfauthor={Léo Martin}}
\pagestyle{empty}
\setlength{\parindent}{0pt}
\setlength{\parskip}{0pt}
\setlength{\emergencystretch}{1em}
\linespread{1.08}
\setlist[itemize]{leftmargin=1.2em,itemsep=3pt,topsep=5pt,parsep=0pt,partopsep=0pt}
\titleformat{\section}{\large\bfseries\color{accent}}{}{0pt}{}[\vspace{2pt}\titlerule]
\titlespacing*{\section}{0pt}{14pt plus 0.7fill}{8pt}
\newcommand{\experience}[3]{\textbf{#1}\hfill{\small\color{muted}#2}\par\textit{#3}\par}
\newcommand{\experiencegap}{\vspace{8pt plus 1fill}}
\begin{document}
{\fontsize{27}{30}\selectfont\bfseries Product Builder Engineer}\par
\vspace{6pt}
{\large\color{accent}Léo Martin}\hfill{\color{muted}Nantes (St-Herblain)}\par
\vspace{7pt}
\href{mailto:lmartin.leo@gmail.com}{lmartin.leo@gmail.com}\quad\textbar\quad\href{https://me.filou.org}{me.filou.org}\quad\textbar\quad\href{https://github.com/OhohLeo}{github.com/OhohLeo}\par
\vspace{8pt}
{\color{muted}Go / Python – IA – Microservices – Cloud – IoT}\par
\vspace{8pt}
Ingénieur passionné par la création de produits de bout en bout, de l'embarqué au cloud, combinant IoT, backend~et~IA.\par

\section{Expériences professionnelles}
\experience{Robeau}{2025 – aujourd'hui}{CTO fondateur – plateforme IoT supervisant + de 2 000 équipements}
\begin{itemize}
\item Conception et orchestration d'une plateforme IoT microservices : Python, PostgreSQL et NATS.
\item Intégration des interfaces React et de solutions IA (Claude, Google Document AI, TimeFS).
\item Déploiement et supervision cloud : GCE, Docker, Grafana, Ansible ; sécurité avec Trivy et ZAP.
\item Déploiement d'une plateforme de gestion des passerelles LoRaWAN.
\item Back-office ERPNext : gestion des stocks et des interventions, mobilapp \& site personnalisé.
\end{itemize}
\experiencegap

\experience{QOS Energy / Univers}{2018 – 2024}{Acteur clé dans la conversion d'un monolithe en microservices}
\begin{itemize}
\item Migration d'un monolithe vers des microservices Go/gRPC sur GKE et Kubernetes.
\item Traitement de flux Kafka (Go, Scala), APIs REST (Go, Ruby/Rails, MySQL, PostgreSQL).
\item Plateforme d'exécution d'algorithmes en Python et Redis.
\end{itemize}
\experiencegap

\experience{Entech}{2017}{Concepteur d'une solution de supervision adaptable et facilement configurable}
\begin{itemize}
\item Supervisions embarquées configurables : Qt/C++, Perl, Raspberry Pi et Modbus.
\item Simulateurs de test en Go/Angular et mise en service chez les clients.
\end{itemize}
\experiencegap

\experience{Ijenko}{2009 – 2016}{Développeur full-stack : embarqué, backend \& interface utilisateur}
\begin{itemize}
\item Plateforme domotique multi-protocolaire en Perl et APIs REST en Go/Perl/MySQL.
\item WebApp mobile Angular/TypeScript ; protocole ZigBee en C et intégration de capteurs tiers.
\end{itemize}
\experiencegap

\section{Compétences}
Architecture · Go · Python · TypeScript / React · PostgreSQL\par
API REST / gRPC · Kafka / NATS · Docker · Git / GitLab-CI\par

\section{Formation}
\textbf{INSA de Rennes}\hfill2009\par
Diplôme d'ingénieur – Systèmes et Réseaux de Communication\par

\section{Langues}
Français (Langue maternelle) · Anglais (980 TOEIC)\par
\end{document}
